% !PRINTED 185
\weberpage{185}

\begin{center}
§. 59.

\gesperrt{Transformation of forms of degree $n$.}
\end{center}

When a form of degree $n$ in $m$ variables, $F(x)$, is transformed by a linear substitution [\S.\,55, (1)] into $\Phi(x')$, then at the same time the linear function of the variables $\xi_1, \xi_2 \ldots \xi_m$

\weq{1}{\overset{i}{\varSigma}\, \xi_i F'(x_i)}

is transformed, by the same substitution applied to the $\xi_i$, into

% [sic] In (2) the print omits the derivative prime on \Phi, although (1), the coefficient argument, and the following identity require \Phi'. The print is preserved.
\weq{2}{\overset{i}{\varSigma}\, \xi_i' \Phi(x_i')}

because both expressions represent the coefficient of the first power of $t$ in the expansion of
\[
F(x_i + t\xi_i) = \Phi(x_i' + t\xi_i')
\]
in powers of $t$ (\S.\,16). Likewise the quadratic function

\weq{3}{\varSigma\, \xi_i \xi_k F''(x_i x_k)}

is transformed into

\weq{4}{\varSigma\, \xi_i' \xi_k' \Phi''(x_i', x_k')}

If we apply the first result to a system of $m$ functions $F_1(x), F_2(x) \ldots F_m(x)$ of arbitrary degrees, and denote by $\varDelta$ the determinant
\[
\begin{vmatrix}
F_1'(x_1),\ F_1'(x_2) \ldots F_1'(x_m) \\
F_2'(x_1),\ F_2'(x_2) \ldots F_2'(x_m) \\
\cdot\ \cdot\ \cdot\ \cdot\ \cdot\ \cdot\ \cdot\ \cdot\ \cdot\ \cdot \\
F_m'(x_1),\ F_m'(x_2) \ldots F_m'(x_m)
\end{vmatrix}\! ,
\]
while we denote by $\varDelta'$ the corresponding determinant formed from the functions $\Phi_1, \Phi_2 \ldots \Phi_m$, then by \S.\,55, (4), if we again denote by $r$ the substitution determinant, we obtain

\weq{5}{\varDelta' = r\varDelta.}

The function $\varDelta$ is called the \gesperrt{functional determinant} of the $m$ functions $F_1, F_2 \ldots F_m$.

Likewise we obtain from the transformation (3), (4), if we denote by $H$ the determinant
%CONT>>
